% File: part2/chapters3/intro.tex

\chapter{TỐI ƯU HÓA, TRIỂN KHAI VÀ VẬN HÀNH (MLOps)}
\label{chap:mlops}

Công việc với một mô hình học máy không kết thúc ở epoch huấn luyện cuối cùng hay ở con số 95\% F1-score trên tập kiểm tra. Mô hình chỉ có giá trị khi được đưa vào một sản phẩm, phục vụ hàng nghìn, thậm chí hàng triệu người dùng, và hoạt động đáng tin cậy 24/7. Phần việc đó gọi là MLOps, vận hành học máy.

Chương này bàn ba vấn đề nằm ngoài khuôn khổ một sổ tay Jupyter: cách phục vụ một mô hình 70 tỷ tham số với độ trễ thấp, cách đóng gói một ứng dụng AI phức tạp để nó chạy được ở bất cứ đâu, và cách giám sát mô hình sau khi triển khai để biết khi nào cần huấn luyện lại.

Các thành phần cơ sở hạ tầng được trình bày gồm cơ sở dữ liệu vector, các thư viện tối ưu hóa mô hình, các framework phục vụ chuyên dụng, và cuối cùng là một vòng đời MLOps hoàn chỉnh cho LLM.